\documentclass[11pt]{article}

\usepackage[margin=1in]{geometry}
\usepackage{times}
\usepackage{amsmath,amssymb}
\usepackage{graphicx}
\usepackage{booktabs}
\usepackage{enumitem}
\usepackage{hyperref}
\usepackage[authoryear,round]{natbib}

\title{
Refusal Evolution over Multiple Turns in Large Language Models\\
\vspace{0.3cm}
\large Research Proposal
}

\author{}

\date{}

\begin{document}

\maketitle

\begin{abstract}

Current work on jailbreak attacks primarily evaluates success rates, while the 
internal computational mechanisms that allow alignment to degrade across long 
conversations remain poorly understood.
This proposal aims to develop mechanistic explanations for the evolution of 
refusal and compliance during multi-turn interactions, we primarily 
focus on features, circuits and universality in multi turn interactions
rather than just proving causality.

\end{abstract}

\section{Introduction}

Prior work has shown that Large Language Models remain aligned across isolated 
single turn prompts, however these results do not extend to multi-turn 
interactions that more accurately mimic real world use cases.
Little is known about the internal representations of refusal in multi-turn
interactions, how it evolves over turns and what recieves attention and how that
affects the alignment. 
We aim to provide a mechanistic understanding on the evolution of refusal
in multi turn Large Language Model interactions rather than just proving 
the existence of correlation between conversational input in multi-turn interactions
and the corresponding effects on alignment. 

\section{Hypothesis}

We hypothesize that multi-turn jailbreaks do not create entirely new 
behaviors but instead gradually suppress or overwrite internal refusal 
representations through identifiable computational pathways.
Furthermore, these mechanisms should exhibit significant similarity across 
related instruction-tuned language models.

\section{Research Questions}

The project investigates the following questions:
\begin{itemize}

    \item
    Which transformer layers exhibit the earliest measurable alignment drift?

    \item
    Are observed changes explained by a small number of interpretable features, 
    circuits, or activation directions?

    \item
    Can causal interventions recover refusal after successful jailbreaks and why?

    \item 
    Does refusal weaken over time or if it is overtaken by competing objectives 
    such as role-playing, instruction following or helpfulness?

    \item 
    How does information from earlier conversation turns propagate through the
    Large Language Model to determine whether refusal depends primarily on recent
    context, conversation history or cumulative conversational effects? 

\end{itemize}

\section{Related Work}

Relevant areas include Alignment drift, Multi-turn jailbreak attacks, 
Refusal directions, Sparse autoencoder features, Activation steering,
Transformer circuit analysis.

\section{Methodology}

\begin{itemize}
    
    \item 
    Use techniques such as Logit Lens \citep{nostalgebraist2020} and 
    Attention Lens \citep{sakarvadia2023attentionlenstoolmechanistically} 
    to probe activations and attention heads. 
    These would be supplemented by Logit Lens Translators 
    \citep{yom-din-etal-2024-jump, belrose2025elicitinglatentpredictionstransformers} 
    and Logit spectroscopy \citep{cancedda-2024-spectral} to ensure accurate 
    unembedding by transforming the intermediate activations to one 
    that matches more accurately to the final layer activations and to discover
    overlooked information by spectral filtering.

    \item 
    Use intervention based ablation \citep{chan2022causal} specifically 
    resampling ablation to prevent the model from going out of distribution.
    Employ activation patching \citep{3600270.3601532}, path patching 
    \citep{wang2022interpretability}, and distributed alignment search 
    \citep{geiger2024findingalignmentsinterpretablecausal} for patching different
    granularity levels and edge attribution patching for automating interventions 
    for localization. 
    The patching effect to evaluate importance of the features is measured by
    Probability or Logit Difference and KL Divergence.

    \item 
    This will be further tested with SAEs \citep{cunningham2023sparse}


\end{itemize}

\subsection{Models}

Potential evaluation models include

\begin{itemize}

\item Llama 3

\item Qwen

\item Gemma

\item Mistral

\end{itemize}

\subsection{Datasets}

Experiments will include

\begin{itemize}

\item MultiBreak

\item HarmBench

\item AdvBench

\item Crescendo conversations

\end{itemize}

\subsection{Analysis Pipeline}

For every conversation:

\begin{enumerate}

\item Run the complete dialogue.

\item Save activations from every transformer block.

\item Measure refusal/compliance scores.

\item Track latent representation evolution.

\item Identify important neurons or SAE features.

\item Perform causal interventions.

\item Measure behavioral recovery.

\end{enumerate}

\section{Evaluation}

Primary evaluation metrics include

\begin{itemize}

\item Jailbreak success rate

\item Refusal score

\item Representation similarity

\item Feature sparsity

\item Causal effect size

\item Layer-wise alignment drift

\end{itemize}

\section{Expected Contributions}

This work aims to contribute

\begin{enumerate}

\item A mechanistic characterization of multi-turn alignment evolution.

\item Benchmarks for layer-wise refusal dynamics.

\item Open-source analysis tools.

\item New hypotheses regarding alignment computation.

\end{enumerate}

\section{Risks}

Potential challenges include

\begin{itemize}

\item Difficulty identifying consistent mechanisms

\item Model-specific behaviors

\item Large computational requirements

\item Ambiguous causal attribution

\end{itemize}

\section{Broader Impact}

Understanding internal alignment mechanisms may improve future safety interventions, enable more interpretable defenses against jailbreaks, and contribute toward principled alignment research.

\bibliographystyle{plainnat}
\bibliography{references}

\end{document}